% Step 1.2: bolus simulation and rotational DSA rendering.
% Numbers from results/step1.2_bolus_summary.json, results/step1.2_render_summary.json.
% Bib keys to add to the .bib: thakur2026punch, liu2025fourdrgs, biguri2016tigre, zha2024r2gaussian.

\paragraph{Contrast transport.}
To avoid an inverse crime, the data are generated with a transport model
that is richer than the plug flow used by FlowSplat. For each case we solve
the area-aware one-dimensional advection--dispersion equation of
PUNCH~\cite{thakur2026punch} at the resting flow, in conservative form
$\partial_t(Ac) + \partial_s(Qc) = \partial_s(AD\,\partial_s c)$, on a
0.1\,mm grid with a second-order limited finite-volume scheme. The
dispersion coefficient is drawn log-uniformly over one decade around
$10^{-6}$\,m$^2$/s~\cite{thakur2026punch}, and the inlet is a gamma-variate
bolus with randomized arrival time and shape, perturbed by 5\,\% smooth
multiplicative noise. In all cases both wash-in and wash-out occur within
the acquisition.

\paragraph{Rotational DSA.}
Projections follow the geometry of 4DRGS~\cite{liu2025fourdrgs}: 133
frames evenly spaced over $198^\circ$, with source-to-object and
source-to-detector distances of 750\,mm and 1200\,mm. We assume a sweep
duration of 4.4\,s, so that each frame is acquired at its own angle and
time. For every frame the lumen,
voxelized at 0.25\,mm with partial volume from the centerline and the
stenosed radius, is filled with the simulated concentration and projected
with TIGRE~\cite{biguri2016tigre} onto a detector with 0.64\,mm pixels;
Poisson and Gaussian noise are added as in
R$^2$-Gaussian~\cite{zha2024r2gaussian}. Training subsets use 30, 60 or
all 133 evenly spaced views, the remaining frames serving as held-out
views. The 200 cases are split by patient into 140 training, 30
validation and 30 test cases.
